\section{Normal cofinal cardinal sequences at singular cardinals}

\import{ordinal.tex}
\import{cardinal/comparison.tex}
\import{cardinal/normal-map.tex}
\import{cardinal/regularity.tex}
\import{recursion/transfinite.tex}

The bounded recursion below constructs a normal cofinal cardinal sequence at
every singular cardinal.  Its values are in fact infinite cardinals, and the
sequence is index-dominating; these are separate properties supplied by the
construction rather than parts of normality.

\begin{definition}\label{cardinal_sequence_envelope}
    $\cardinalSequenceEnvelope{\kappa}{d}
        =\cofinalEnvelope{d}{\cofinality{\kappa}}{\kappa}$.
\end{definition}

\begin{definition}\label{cardinal_sequence_next}
    $\cardinalSequenceNext{q}{\lambda}{\beta}{g}
        =\cardSucc{\card{(\lambda\union q(\beta))
            \union\unions{\ran{g}}}}$.
\end{definition}

\begin{definition}\label{cardinal_sequence_limit_stage}
    $\cardinalSequenceLimitStage{\beta}$ iff
        $\beta$ is a limit ordinal.
\end{definition}

\begin{definition}\label{cardinal_sequence_stage_property}
    $\cardinalSequenceStageProperty{c}{\beta}$ iff
        $c(\beta)$ is an infinite cardinal and
        for every $j\in\beta$ we have $c(j)\in c(\beta)$.
\end{definition}

\begin{definition}\label{cardinal_sequence_step}
    $\cardinalSequenceStep{q}{\lambda}{\kappa}
        =\{v\in((\lambda\times\pow{\lambda\times\kappa})\times\kappa)
            \mid
            (\cardinalSequenceLimitStage{\fst{\fst{v}}}\land
             \unions{\ran{\snd{\fst{v}}}}\in\kappa\land
             \snd{v}=\unions{\ran{\snd{\fst{v}}}})\lor
            (\cardinalSequenceLimitStage{\fst{\fst{v}}}\land
             \unions{\ran{\snd{\fst{v}}}}\notin\kappa\land
             \snd{v}=\emptyset)\lor
            (\lnot\cardinalSequenceLimitStage{\fst{\fst{v}}}\land
             \cardinalSequenceNext{q}{\lambda}{\fst{\fst{v}}}
                {\snd{\fst{v}}}\in\kappa\land
             \snd{v}=\cardinalSequenceNext{q}{\lambda}{\fst{\fst{v}}}
                {\snd{\fst{v}}})\lor
            (\lnot\cardinalSequenceLimitStage{\fst{\fst{v}}}\land
             \cardinalSequenceNext{q}{\lambda}{\fst{\fst{v}}}
                {\snd{\fst{v}}}\notin\kappa\land
             \snd{v}=\emptyset)
        \}$.
\end{definition}

\begin{definition}\label{cardinal_sequence}
    $\cardinalSequence{\kappa}{d}
        =\ordinalRecursion{\cofinality{\kappa}}{\kappa}
            {\cardinalSequenceStep{\cardinalSequenceEnvelope{\kappa}{d}}
                {\cofinality{\kappa}}{\kappa}}$.
\end{definition}

\begin{proposition}\label{cardinal_sequence_step_type}
    Let $\kappa$ be an infinite cardinal.
    Then $\cardinalSequenceStep{q}{\lambda}{\kappa}
        \in\funs{\lambda\times\pow{\lambda\times\kappa}}{\kappa}$.
\end{proposition}
\begin{proof}
    Let $D=\lambda\times\pow{\lambda\times\kappa}$.
    Let $H=\cardinalSequenceStep{q}{\lambda}{\kappa}$.
    $H\subseteq D\times\kappa$
        by \auto{cardinal_sequence_step,subseteq}.
    Hence $H$ is a relation
        by \auto{relation,subseteq,times_elem_is_tuple}.
    Show that $H$ is right-unique.
    \begin{subproof}
        Show that for all $u,b,c$ such that $(u,b),(u,c)\in H$
            we have $b=c$.
        \begin{subproof}
            Fix $u,b,c$.
            Assume $(u,b),(u,c)\in H$.
            Let $\beta=\fst{u}$.
            Let $g=\snd{u}$.
            \begin{byCase}
                \caseOf{$\cardinalSequenceLimitStage{\beta}$.}
                    \begin{byCase}
                        \caseOf{$\unions{\ran{g}}\in\kappa$.}
                            $b=\unions{\ran{g}}$ by
                                \auto{cardinal_sequence_step,fst_eq,snd_eq,pair_eq_iff}.
                            $c=\unions{\ran{g}}$ by
                                \auto{cardinal_sequence_step,fst_eq,snd_eq,pair_eq_iff}.
                        \caseOf{$\unions{\ran{g}}\notin\kappa$.}
                            $b=\emptyset$ by
                                \auto{cardinal_sequence_step,fst_eq,snd_eq,pair_eq_iff}.
                            $c=\emptyset$ by
                                \auto{cardinal_sequence_step,fst_eq,snd_eq,pair_eq_iff}.
                    \end{byCase}
                \caseOf{$\lnot\cardinalSequenceLimitStage{\beta}$.}
                    \begin{byCase}
                        \caseOf{$\cardinalSequenceNext{q}{\lambda}{\beta}{g}
                            \in\kappa$.}
                            $b=\cardinalSequenceNext{q}{\lambda}{\beta}{g}$ by
                                \auto{cardinal_sequence_step,fst_eq,snd_eq,pair_eq_iff}.
                            $c=\cardinalSequenceNext{q}{\lambda}{\beta}{g}$ by
                                \auto{cardinal_sequence_step,fst_eq,snd_eq,pair_eq_iff}.
                        \caseOf{$\cardinalSequenceNext{q}{\lambda}{\beta}{g}
                            \notin\kappa$.}
                            $b=\emptyset$ by
                                \auto{cardinal_sequence_step,fst_eq,snd_eq,pair_eq_iff}.
                            $c=\emptyset$ by
                                \auto{cardinal_sequence_step,fst_eq,snd_eq,pair_eq_iff}.
                    \end{byCase}
            \end{byCase}
        \end{subproof}
        Follows by \obv{rightunique}.
    \end{subproof}
    Show that $\dom{H}=D$.
    \begin{subproof}
        $\dom{H}\subseteq D$
            by \auto{subseteq,dom_iff,cardinal_sequence_step,times_tuple_elim,pair_eq_iff}.
        Show that every element of $D$ is an element of $\dom{H}$.
        \begin{subproof}
            Fix $u\in D$.
            Let $\beta=\fst{u}$.
            Let $g=\snd{u}$.
            $u=(\beta,g)$ by \auto{times_elem_is_tuple,pair_eq_pair_of_proj}.
            $\emptyset\in\kappa$
                by \obv{zero_one_in_infinite_cardinal}.
            \begin{byCase}
                \caseOf{$\cardinalSequenceLimitStage{\beta}$.}
                    \begin{byCase}
                        \caseOf{$\unions{\ran{g}}\in\kappa$.}
                            $(u,\unions{\ran{g}})\in H$
                                by \auto{cardinal_sequence_step,fst_eq,snd_eq,times_tuple_intro}.
                            Follows by \obv{dom_intro}.
                        \caseOf{$\unions{\ran{g}}\notin\kappa$.}
                            $(u,\emptyset)\in H$
                                by \auto{cardinal_sequence_step,fst_eq,snd_eq,times_tuple_intro}.
                            Follows by \obv{dom_intro}.
                    \end{byCase}
                \caseOf{$\lnot\cardinalSequenceLimitStage{\beta}$.}
                    \begin{byCase}
                        \caseOf{$\cardinalSequenceNext{q}{\lambda}{\beta}{g}
                            \in\kappa$.}
                            $(u,\cardinalSequenceNext{q}{\lambda}{\beta}{g})
                                \in H$
                                by \auto{cardinal_sequence_step,fst_eq,snd_eq,times_tuple_intro}.
                            Follows by \obv{dom_intro}.
                        \caseOf{$\cardinalSequenceNext{q}{\lambda}{\beta}{g}
                            \notin\kappa$.}
                            $(u,\emptyset)\in H$
                                by \auto{cardinal_sequence_step,fst_eq,snd_eq,times_tuple_intro}.
                            Follows by \obv{dom_intro}.
                    \end{byCase}
            \end{byCase}
        \end{subproof}
        Hence $D\subseteq\dom{H}$ by \obv{subseteq}.
        Follows by \obv{subseteq_antisymmetric}.
    \end{subproof}
    For every $u\in D$ we have $H(u)\in\kappa$
        by \auto{function_apply_elim,subseteq,times_tuple_elim}.
    Hence $H$ is a function to $\kappa$.
    Follows by \auto{funs_intro}.
\end{proof}

\begin{proposition}\label{cardinal_sequence_step_limit}
    Let $\kappa$ be an infinite cardinal.
    Let $\beta\in\lambda$.
    Suppose $\beta$ is a limit ordinal.
    Let $g\in\pow{\lambda\times\kappa}$.
    Suppose $\unions{\ran{g}}\in\kappa$.
    Then $(\cardinalSequenceStep{q}{\lambda}{\kappa})((\beta,g))
        =\unions{\ran{g}}$.
\end{proposition}
\begin{proof}
    $(\beta,g)\in\lambda\times\pow{\lambda\times\kappa}$
        by \obv{times_tuple_intro}.
    $((\beta,g),\unions{\ran{g}})
        \in\cardinalSequenceStep{q}{\lambda}{\kappa}$
        by \auto{cardinal_sequence_step,fst_eq,snd_eq,cardinal_sequence_limit_stage,times_tuple_intro}.
    $\cardinalSequenceStep{q}{\lambda}{\kappa}$ is a function
        by \auto{cardinal_sequence_step_type,funs_is_function}.
    Follows by \obv{function_apply_intro}.
\end{proof}

\begin{proposition}\label{cardinal_sequence_step_next}
    Let $\kappa$ be an infinite cardinal.
    Let $\beta\in\lambda$.
    Suppose $\beta$ is not a limit ordinal.
    Let $g\in\pow{\lambda\times\kappa}$.
    Suppose $\cardinalSequenceNext{q}{\lambda}{\beta}{g}\in\kappa$.
    Then $(\cardinalSequenceStep{q}{\lambda}{\kappa})((\beta,g))
        =\cardinalSequenceNext{q}{\lambda}{\beta}{g}$.
\end{proposition}
\begin{proof}
    $(\beta,g)\in\lambda\times\pow{\lambda\times\kappa}$
        by \obv{times_tuple_intro}.
    $((\beta,g),\cardinalSequenceNext{q}{\lambda}{\beta}{g})
        \in\cardinalSequenceStep{q}{\lambda}{\kappa}$
        by \auto{cardinal_sequence_step,fst_eq,snd_eq,cardinal_sequence_limit_stage,times_tuple_intro}.
    $\cardinalSequenceStep{q}{\lambda}{\kappa}$ is a function
        by \auto{cardinal_sequence_step_type,funs_is_function}.
    Follows by \obv{function_apply_intro}.
\end{proof}

\begin{proposition}\label{cardinal_sequence_type}
    Let $\kappa$ be a singular cardinal.
    Let $\lambda=\cofinality{\kappa}$.
    Suppose $\cofinalMap{d}{\lambda}{\kappa}$.
    Then $\cardinalSequence{\kappa}{d}\in\funs{\lambda}{\kappa}$.
\end{proposition}
\begin{proof}
    $\lambda$ is a cardinal by \obv{cofinality_is_cardinal}.
    Hence $\lambda$ is an ordinal by \auto{cardinal,initial_ordinal}.
    $\cardinalSequenceStep{\cardinalSequenceEnvelope{\kappa}{d}}
        {\lambda}{\kappa}
        \in\funs{\lambda\times\pow{\lambda\times\kappa}}{\kappa}$
        by \obv{cardinal_sequence_step_type}.
    Follows by
        \auto{cardinal_sequence,ordinal_recursion_type}.
\end{proof}

\begin{proposition}\label{cardinal_sequence_equation}
    Let $\kappa$ be a singular cardinal.
    Let $\lambda=\cofinality{\kappa}$.
    Suppose $\cofinalMap{d}{\lambda}{\kappa}$.
    Let $\beta\in\lambda$.
    Let $c=\cardinalSequence{\kappa}{d}$.
    Then $c(\beta)=
        (\cardinalSequenceStep{\cardinalSequenceEnvelope{\kappa}{d}}
            {\lambda}{\kappa})((\beta,\restrl{c}{\beta}))$.
\end{proposition}
\begin{proof}
    $\lambda$ is a cardinal by \obv{cofinality_is_cardinal}.
    Hence $\lambda$ is an ordinal by \auto{cardinal,initial_ordinal}.
    $\cardinalSequenceStep{\cardinalSequenceEnvelope{\kappa}{d}}
        {\lambda}{\kappa}
        \in\funs{\lambda\times\pow{\lambda\times\kappa}}{\kappa}$
        by \obv{cardinal_sequence_step_type}.
    Follows by
        \auto{cardinal_sequence,ordinal_recursion_equation}.
\end{proof}

\begin{proposition}\label{cardinal_sequence_envelope_type}
    Let $\kappa$ be a singular cardinal.
    Let $\lambda=\cofinality{\kappa}$.
    Suppose $\cofinalMap{d}{\lambda}{\kappa}$.
    Then $\cardinalSequenceEnvelope{\kappa}{d}
        \in\funs{\lambda}{\kappa}$.
\end{proposition}
\begin{proof}
    Hence $\kappa$ is a limit ordinal
        by \obv{infinite_cardinal_is_limit_ordinal}.
    Hence $\kappa$ is an ordinal by \obv{limit_ordinal}.
    Follows by
        \auto{cardinal_sequence_envelope,cofinal_envelope_type}.
\end{proof}

\begin{proposition}\label{cardinal_sequence_envelope_dominates}
    Let $\kappa$ be a singular cardinal.
    Let $\lambda=\cofinality{\kappa}$.
    Suppose $\cofinalMap{d}{\lambda}{\kappa}$.
    Let $\beta\in\lambda$.
    Let $j\in\suc{\beta}$.
    Then $d(j)\in(\cardinalSequenceEnvelope{\kappa}{d})(\beta)$.
\end{proposition}
\begin{proof}
    Hence $\kappa$ is a limit ordinal
        by \obv{infinite_cardinal_is_limit_ordinal}.
    Follows by
        \auto{cardinal_sequence_envelope,cofinal_envelope_dominates_prefix}.
\end{proof}

\begin{proposition}\label{union_of_cardinals_is_cardinal}
    Suppose every element of $C$ is a cardinal.
    Then $\unions{C}$ is a cardinal.
\end{proposition}
\begin{proof}
    Let $U=\unions{C}$.
    Every element of $C$ is an ordinal
        by \auto{cardinal,initial_ordinal}.
    Hence $U$ is an ordinal by \obv{union_of_ordinals_is_ordinal}.
    Show that for every $\alpha\in U$ we have $\lnot(\alpha\approx U)$.
    \begin{subproof}
        Fix $\alpha\in U$.
        Take $\kappa\in C$ such that $\alpha\in\kappa$
            by \obv{unions_iff}.
        Hence $\kappa$ is an initial ordinal by \obv{cardinal}.
        Suppose $\alpha\approx U$.
        $\kappa\subseteq U$ by \auto{subseteq,unions_iff}.
        Hence $\cardLe{\kappa}{U}$ by \obv{subseteq_cardinal_le}.
        Hence $\cardLe{U}{\alpha}$ by \auto{equinum_cardinal_le}.
        Hence $\cardLe{\kappa}{\alpha}$
            by \obv{cardinal_le_transitive}.
        $\alpha$ and $\kappa$ are ordinals
            by \obv{initial_ordinal,elem_of_ordinal_is_ordinal}.
        Hence $\alpha\subseteq\kappa$
            by \obv{ordinal_elem_implies_subseteq}.
        Hence $\cardLe{\alpha}{\kappa}$
            by \obv{subseteq_cardinal_le}.
        Hence $\kappa\approx\alpha$
            by \obv{cardinal_le_antisymmetric_up_to_equinum}.
        Hence $\alpha\approx\kappa$ by \obv{equinum_sym}.
        Contradiction by \obv{initial_ordinal}.
    \end{subproof}
    Hence $U$ is an initial ordinal by \obv{initial_ordinal}.
    Follows by \obv{cardinal}.
\end{proof}

\begin{proposition}\label{successor_cardinal_below_limit_cardinal}
    Let $\kappa$ be a limit cardinal.
    Let $\mu\in\kappa$.
    Suppose $\mu$ is a cardinal.
    Then $\cardSucc{\mu}\in\kappa$.
\end{proposition}
\begin{proof}
    $\card{\mu}=\mu$ and $\card{\kappa}=\kappa$
        by \obv{card_of_cardinal}.
    Hence $\cardLt{\mu}{\kappa}$
        by \auto{card_elem_iff_cardinal_lt}.
    Hence $\cardSucc{\mu}\subseteq\kappa$
        by \obv{cardinal_successor_minimal}.
    $\cardSucc{\mu}$ and $\kappa$ are ordinals
        by \auto{cardinal_successor_is_cardinal,cardinal,initial_ordinal}.
    $\cardSucc{\mu}\neq\kappa$ by \auto{limit_cardinal}.
    Follows by \obv{ordinal_subseteq_implies_elem_or_eq}.
\end{proof}

\begin{proposition}\label{cardinal_sequence_next_below}
    Let $\kappa$ be a singular cardinal.
    Let $\lambda=\cofinality{\kappa}$.
    Let $q\in\funs{\lambda}{\kappa}$.
    Let $\beta\in\lambda$.
    Let $g\in\funs{\beta}{\kappa}$.
    Then $\cardinalSequenceNext{q}{\lambda}{\beta}{g}\in\kappa$.
\end{proposition}
\begin{proof}
    $\kappa$ is a limit cardinal
        by \obv{singular_cardinal_is_limit_cardinal}.
    Hence $\kappa$ is a limit ordinal
        by \obv{infinite_cardinal_is_limit_ordinal}.
    Hence $\kappa$ is an ordinal by \obv{limit_ordinal}.
    $\lambda\in\kappa$
        by \auto{cofinality_belongs_to_singular_cardinal}.
    $q(\beta)\in\kappa$ by \obv{funs_type_apply}.
    $\unions{\ran{g}}\in\kappa$
        by \obv{small_domain_range_union_below}.
    $(\lambda\union q(\beta))\union\unions{\ran{g}}\in\kappa$
        by \obv{three_ordinal_union_below}.
    Let $B=(\lambda\union q(\beta))\union\unions{\ran{g}}$.
    $B\in\kappa$.
    $\card{B}\in\kappa$ by \obv{card_of_ordinal_below}.
    $\card{B}$ is a cardinal by \obv{card_is_cardinal}.
    Hence $\cardSucc{\card{B}}\in\kappa$
        by \obv{successor_cardinal_below_limit_cardinal}.
    Follows by \auto{cardinal_sequence_next}.
\end{proof}

\begin{proposition}\label{cardinal_sequence_restriction_type}
    Let $\kappa$ be a singular cardinal.
    Let $\lambda=\cofinality{\kappa}$.
    Suppose $\cofinalMap{d}{\lambda}{\kappa}$.
    Let $\beta\in\lambda$.
    Let $c=\cardinalSequence{\kappa}{d}$.
    Then $\restrl{c}{\beta}\in\funs{\beta}{\kappa}$.
\end{proposition}
\begin{proof}
    $c\in\funs{\lambda}{\kappa}$ by \auto{cardinal_sequence_type}.
    $\lambda$ is an ordinal
        by \auto{cofinality_is_cardinal,cardinal,initial_ordinal}.
    Hence $\beta$ is an ordinal and $\beta\subseteq\lambda$
        by \obv{elem_of_ordinal_is_ordinal,ordinal_elem_implies_subseteq}.
    Follows by \obv{funs_restrl}.
\end{proof}

\begin{proposition}\label{cardinal_sequence_at_limit}
    Let $\kappa$ be a singular cardinal.
    Let $\lambda=\cofinality{\kappa}$.
    Suppose $\cofinalMap{d}{\lambda}{\kappa}$.
    Let $\beta\in\lambda$.
    Suppose $\beta$ is a limit ordinal.
    Let $c=\cardinalSequence{\kappa}{d}$.
    Then $c(\beta)=\unions{\ran{\restrl{c}{\beta}}}$.
\end{proposition}
\begin{proof}
    Let $g=\restrl{c}{\beta}$.
    $g\in\funs{\beta}{\kappa}$
        by \auto{cardinal_sequence_restriction_type}.
    Hence $g\in\rels{\beta}{\kappa}$
        by \obv{funs_subseteq_rels,subseteq}.
    Hence $g\subseteq\beta\times\kappa$ by \obv{rels_elim}.
    $\lambda$ is an ordinal
        by \auto{cofinality_is_cardinal,cardinal,initial_ordinal}.
    Hence $\beta\subseteq\lambda$
        by \obv{elem_of_ordinal_is_ordinal,ordinal_elem_implies_subseteq}.
    Hence $\beta\times\kappa\subseteq\lambda\times\kappa$
        by \obv{times_subseteq_left}.
    Hence $g\subseteq\lambda\times\kappa$
        by \obv{subseteq_transitive}.
    Hence $g\in\pow{\lambda\times\kappa}$ by \obv{powerset_intro}.
    $\kappa$ is an infinite cardinal by \obv{singular_cardinal}.
    Hence $\kappa$ is a limit ordinal
        by \obv{infinite_cardinal_is_limit_ordinal}.
    $\unions{\ran{g}}\in\kappa$
        by \obv{small_domain_range_union_below}.
    $c(\beta)=
        (\cardinalSequenceStep{\cardinalSequenceEnvelope{\kappa}{d}}
            {\lambda}{\kappa})((\beta,g))$
        by \auto{cardinal_sequence_equation}.
    Hence $c(\beta)=\unions{\ran{g}}$
        by \auto{cardinal_sequence_step_limit}.
\end{proof}

\begin{proposition}\label{cardinal_sequence_at_nonlimit}
    Let $\kappa$ be a singular cardinal.
    Let $\lambda=\cofinality{\kappa}$.
    Suppose $\cofinalMap{d}{\lambda}{\kappa}$.
    Let $\beta\in\lambda$.
    Suppose $\beta$ is not a limit ordinal.
    Let $c=\cardinalSequence{\kappa}{d}$.
    Let $q=\cardinalSequenceEnvelope{\kappa}{d}$.
    Then $c(\beta)=
        \cardinalSequenceNext{q}{\lambda}{\beta}{\restrl{c}{\beta}}$.
\end{proposition}
\begin{proof}
    Let $g=\restrl{c}{\beta}$.
    $g\in\funs{\beta}{\kappa}$
        by \auto{cardinal_sequence_restriction_type}.
    Hence $g\in\rels{\beta}{\kappa}$
        by \obv{funs_subseteq_rels,subseteq}.
    Hence $g\subseteq\beta\times\kappa$ by \obv{rels_elim}.
    $\lambda$ is an ordinal
        by \auto{cofinality_is_cardinal,cardinal,initial_ordinal}.
    Hence $\beta\subseteq\lambda$
        by \obv{elem_of_ordinal_is_ordinal,ordinal_elem_implies_subseteq}.
    Hence $\beta\times\kappa\subseteq\lambda\times\kappa$
        by \obv{times_subseteq_left}.
    Hence $g\subseteq\lambda\times\kappa$
        by \obv{subseteq_transitive}.
    Hence $g\in\pow{\lambda\times\kappa}$ by \obv{powerset_intro}.
    $q\in\funs{\lambda}{\kappa}$
        by \auto{cardinal_sequence_envelope_type}.
    $\cardinalSequenceNext{q}{\lambda}{\beta}{g}\in\kappa$
        by \obv{cardinal_sequence_next_below}.
    \begin{align*}
        c(\beta)
            &= (\cardinalSequenceStep{q}{\lambda}{\kappa})((\beta,g))
                \explanation{by \auto{cardinal_sequence_equation}}
            \\
            &= \cardinalSequenceNext{q}{\lambda}{\beta}{g}
                \explanation{by \obv{cardinal_sequence_step_next}}
            \\
            &= \cardinalSequenceNext{q}{\lambda}{\beta}{\restrl{c}{\beta}}
    \end{align*}
\end{proof}

\begin{proposition}\label{cardinal_sequence_next_properties}
    Let $\kappa$ be a singular cardinal.
    Let $\lambda=\cofinality{\kappa}$.
    Let $q\in\funs{\lambda}{\kappa}$.
    Let $\beta\in\lambda$.
    Let $g\in\funs{\beta}{\kappa}$.
    Let $h=\cardinalSequenceNext{q}{\lambda}{\beta}{g}$.
    Then $h$ is an infinite cardinal and $\lambda\in h$ and
        $q(\beta)\in h$ and
        for every $j\in\beta$ such that $g(j)$ is a cardinal
            we have $g(j)\in h$.
\end{proposition}
\begin{proof}
    $\kappa$ is a limit cardinal
        by \obv{singular_cardinal_is_limit_cardinal}.
    Hence $\kappa$ is a limit ordinal
        by \obv{infinite_cardinal_is_limit_ordinal}.
    Hence $\kappa$ is an ordinal by \obv{limit_ordinal}.
    $\lambda\in\kappa$
        by \auto{cofinality_belongs_to_singular_cardinal}.
    $q(\beta)\in\kappa$ by \obv{funs_type_apply}.
    Let $U=\unions{\ran{g}}$.
    $U\in\kappa$ by \auto{small_domain_range_union_below}.
    Let $B=(\lambda\union q(\beta))\union U$.
    $(\lambda\union q(\beta))\union U\in\kappa$
        by \obv{three_ordinal_union_below}.
    Hence $B\in\kappa$.
    $\card{B}$ is a cardinal by \obv{card_is_cardinal}.
    Let $H=\cardSucc{\card{B}}$.
    $\cardSucc{\card{B}}$ is a cardinal
        by \obv{cardinal_successor_is_cardinal}.
    $\lambda\subseteq B$
        by \auto{union_upper_left,subseteq_transitive}.
    Hence $\cardLe{\lambda}{B}$ by \obv{subseteq_cardinal_le}.
    Hence $\card{\lambda}\subseteq\card{B}$
        by \obv{card_subseteq_iff_cardinal_le}.
    $\lambda$ is a cardinal by \obv{cofinality_is_cardinal}.
    Hence $\lambda$ is an ordinal and $\card{\lambda}=\lambda$
        by \obv{cardinal_is_ordinal,card_of_cardinal}.
    Hence $\lambda\in H$
        by \obv{ordinal_below_cardinal_successor_iff}.
    $q(\beta)\subseteq B$
        by \auto{union_upper_right,union_upper_left,subseteq_transitive}.
    Hence $\cardLe{q(\beta)}{B}$ by \obv{subseteq_cardinal_le}.
    Hence $\card{q(\beta)}\subseteq\card{B}$
        by \obv{card_subseteq_iff_cardinal_le}.
    $q(\beta)$ is an ordinal by \obv{elem_of_ordinal_is_ordinal}.
    Hence $q(\beta)\in H$
        by \obv{ordinal_below_cardinal_successor_iff}.
    $\lambda$ is an infinite cardinal
        by \obv{cofinality_of_limit_is_infinite_cardinal}.
    Show that $H$ is infinite.
    \begin{subproof}
        Suppose not.
        Hence $H$ is finite.
        $H$ is an ordinal by \auto{cardinal,initial_ordinal}.
        Hence $\lambda\subseteq H$
            by \obv{ordinal_elem_implies_subseteq}.
        Hence $\lambda$ is finite by \obv{subseteq_finite_is_finite}.
        Contradiction.
    \end{subproof}
    Show that for every $j\in\beta$ such that $g(j)$ is a cardinal
        we have $g(j)\in H$.
    \begin{subproof}
        Fix $j$.
        Assume $j\in\beta$ and $g(j)$ is a cardinal.
        $g$ is a function and $\dom{g}=\beta$
            by \obv{funs_is_function,funs}.
        Hence $g(j)\in\ran{g}$ by \auto{fun_ran_iff}.
        Hence $g(j)\subseteq U$ by \auto{subseteq,unions_iff}.
        $U\subseteq B$ by \obv{union_upper_right}.
        Hence $g(j)\subseteq B$ by \obv{subseteq_transitive}.
        Hence $\cardLe{g(j)}{B}$ by \obv{subseteq_cardinal_le}.
        Hence $\card{g(j)}\subseteq\card{B}$
            by \obv{card_subseteq_iff_cardinal_le}.
        $g(j)$ is an ordinal by \auto{cardinal,initial_ordinal}.
        Follows by \obv{ordinal_below_cardinal_successor_iff}.
    \end{subproof}
    $h=H$ by \auto{cardinal_sequence_next}.
\end{proof}

\begin{proposition}\label{cardinal_sequence_stage_step}
    Let $\kappa$ be a singular cardinal.
    Let $\lambda=\cofinality{\kappa}$.
    Suppose $\cofinalMap{d}{\lambda}{\kappa}$.
    Let $c=\cardinalSequence{\kappa}{d}$.
    Let $\beta\in\lambda$.
    Suppose for every $\gamma\in\beta$ we have
        $\cardinalSequenceStageProperty{c}{\gamma}$.
    Then $\cardinalSequenceStageProperty{c}{\beta}$.
\end{proposition}
\begin{proof}
    $c\in\funs{\lambda}{\kappa}$ by \auto{cardinal_sequence_type}.
    Hence $c$ is a function and $\dom{c}=\lambda$
        by \obv{funs_is_function,funs}.
    $\lambda$ is an ordinal
        by \auto{cofinality_is_cardinal,cardinal,initial_ordinal}.
    Hence $\beta$ is an ordinal and $\beta\subseteq\lambda$
        by \obv{elem_of_ordinal_is_ordinal,ordinal_elem_implies_subseteq}.
    Let $g=\restrl{c}{\beta}$.
    $g\in\funs{\beta}{\kappa}$
        by \auto{cardinal_sequence_restriction_type}.
    Hence $g$ is a function and $\dom{g}=\beta$
        by \obv{funs_is_function,funs}.
    For every $j\in\beta$ we have $g(j)=c(j)$
        by \auto{restrl_of_function_apply}.
    \begin{byCase}
        \caseOf{$\beta$ is a limit ordinal.}
            Let $U=\unions{\ran{g}}$.
            $c(\beta)=U$ by \auto{cardinal_sequence_at_limit}.
            Show that every element of $\ran{g}$ is a cardinal.
            \begin{subproof}
                Fix $y\in\ran{g}$.
                Take $\gamma\in\dom{g}$ such that $g(\gamma)=y$
                    by \obv{fun_ran_iff}.
                Hence $\gamma\in\beta$.
                Hence $\cardinalSequenceStageProperty{c}{\gamma}$.
                Hence $c(\gamma)$ is an infinite cardinal
                    by \obv{cardinal_sequence_stage_property}.
                $g(\gamma)=c(\gamma)$.
                Hence $y$ is an infinite cardinal.
            \end{subproof}
            Hence $U$ is a cardinal
                by \obv{union_of_cardinals_is_cardinal}.
            $\emptyset\in\beta$
                by \obv{limit_ordinal}.
            Hence $\cardinalSequenceStageProperty{c}{\emptyset}$.
            $g(\emptyset)=c(\emptyset)$.
            $g(\emptyset)\in\ran{g}$ by \auto{fun_ran_iff}.
            Hence $c(\emptyset)\in\ran{g}$.
            Hence $c(\emptyset)\subseteq U$
                by \auto{subseteq,unions_iff}.
            Show that $U$ is infinite.
            \begin{subproof}
                Suppose not.
                Hence $U$ is finite.
                Hence $c(\emptyset)$ is finite
                    by \obv{subseteq_finite_is_finite}.
                Contradiction.
            \end{subproof}
            Hence $c(\beta)$ is an infinite cardinal.
            Show that for every $j\in\beta$ we have $c(j)\in c(\beta)$.
            \begin{subproof}
                Fix $j\in\beta$.
                $\suc{j}\in\beta$ by \obv{suc_elem_limitordinal}.
                Hence $\cardinalSequenceStageProperty{c}{\suc{j}}$.
                $j\in\suc{j}$ by \obv{suc_intro_self}.
                Hence $c(j)\in c(\suc{j})$
                    by \obv{cardinal_sequence_stage_property}.
                $g(\suc{j})=c(\suc{j})$.
                $g(\suc{j})\in\ran{g}$ by \auto{fun_ran_iff}.
                Hence $c(\suc{j})\in\ran{g}$.
                Hence $c(j)\in U$ by \auto{unions_iff}.
            \end{subproof}
            Follows by \obv{cardinal_sequence_stage_property}.
        \caseOf{$\beta$ is not a limit ordinal.}
            Let $q=\cardinalSequenceEnvelope{\kappa}{d}$.
            $q\in\funs{\lambda}{\kappa}$
                by \auto{cardinal_sequence_envelope_type}.
            Let $h=\cardinalSequenceNext{q}{\lambda}{\beta}{g}$.
            $h$ is an infinite cardinal and $\lambda\in h$ and
                for every $j\in\beta$ such that $g(j)$ is a cardinal
                    we have $g(j)\in h$
                by \auto{cardinal_sequence_next_properties}.
            $c(\beta)=h$ by \auto{cardinal_sequence_at_nonlimit}.
            Hence $c(\beta)$ is an infinite cardinal by \auto.
            Show that for every $j\in\beta$ we have $c(j)\in c(\beta)$.
            \begin{subproof}
                Fix $j\in\beta$.
                Hence $\cardinalSequenceStageProperty{c}{j}$.
                Hence $c(j)$ is an infinite cardinal
                    by \obv{cardinal_sequence_stage_property}.
                $g(j)=c(j)$.
                Hence $g(j)$ is a cardinal by \auto.
                Hence $g(j)\in h$.
            \end{subproof}
            Follows by \obv{cardinal_sequence_stage_property}.
    \end{byCase}
\end{proof}

\begin{theorem}\label{cardinal_sequence_stage_properties}
    Let $\kappa$ be a singular cardinal.
    Let $\lambda=\cofinality{\kappa}$.
    Suppose $\cofinalMap{d}{\lambda}{\kappa}$.
    Let $c=\cardinalSequence{\kappa}{d}$.
    Then for every $\beta\in\lambda$ we have
        $\cardinalSequenceStageProperty{c}{\beta}$.
\end{theorem}
\begin{proof}[Proof by \in-induction on $\beta$]
    Assume $\beta\in\lambda$.
    $\lambda$ is an ordinal
        by \auto{cofinality_is_cardinal,cardinal,initial_ordinal}.
    For every $\gamma\in\beta$ we have $\gamma\in\lambda$
        by \auto{ordinal_transitivity}.
    Hence for every $\gamma\in\beta$ we have
        $\cardinalSequenceStageProperty{c}{\gamma}$.
    Follows by \obv{cardinal_sequence_stage_step}.
\end{proof}

\begin{proposition}\label{cardinal_sequence_contains_cofinality_at_zero}
    Let $\kappa$ be a singular cardinal.
    Let $\lambda=\cofinality{\kappa}$.
    Suppose $\cofinalMap{d}{\lambda}{\kappa}$.
    Let $c=\cardinalSequence{\kappa}{d}$.
    Then $\lambda\in c(\emptyset)$.
\end{proposition}
\begin{proof}
    Hence $\kappa$ is a limit ordinal
        by \obv{infinite_cardinal_is_limit_ordinal}.
    $\lambda$ is an infinite cardinal
        by \obv{cofinality_of_limit_is_infinite_cardinal}.
    Hence $\lambda$ is a limit ordinal
        by \obv{infinite_cardinal_is_limit_ordinal}.
    Hence $\emptyset\in\lambda$ by \obv{limit_ordinal}.
    $\emptyset$ is not a limit ordinal by \obv{zero_not_limitordinal}.
    Let $q=\cardinalSequenceEnvelope{\kappa}{d}$.
    $q\in\funs{\lambda}{\kappa}$
        by \auto{cardinal_sequence_envelope_type}.
    Let $g=\restrl{c}{\emptyset}$.
    $g\in\funs{\emptyset}{\kappa}$
        by \auto{cardinal_sequence_restriction_type}.
    Let $h=\cardinalSequenceNext{q}{\lambda}{\emptyset}{g}$.
    $h$ is an infinite cardinal and $\lambda\in h$
        by \auto{cardinal_sequence_next_properties}.
    $c(\emptyset)=h$ by \auto{cardinal_sequence_at_nonlimit}.
\end{proof}

\begin{lemma}\label{nonlimit_index_above}
    Let $\lambda$ be a limit ordinal.
    Let $j\in\lambda$.
    Then there exists $i\in\lambda$ such that
        $i$ is not a limit ordinal and $j\in\suc{i}$.
\end{lemma}
\begin{proof}
    \begin{byCase}
        \caseOf{$j$ is a limit ordinal.}
            Let $i=\suc{j}$.
            $i\in\lambda$ by \auto{suc_elem_limitordinal}.
            $\suc{j}$ is not a limit ordinal
                by \obv{successor_ordinal_is_not_limit_ordinal}.
            Hence $i$ is not a limit ordinal.
            $j\in i$ by \obv{suc_intro_self}.
            Hence $j\in\suc{i}$ by \auto{suc_iff}.
            Follows by \auto.
        \caseOf{$j$ is not a limit ordinal.}
            $j\in\suc{j}$ by \obv{suc_intro_self}.
            Follows by \auto.
    \end{byCase}
\end{proof}

\begin{proposition}\label{cardinal_sequence_dominates_envelope_at_nonlimit}
    Let $\kappa$ be a singular cardinal.
    Let $\lambda=\cofinality{\kappa}$.
    Suppose $\cofinalMap{d}{\lambda}{\kappa}$.
    Let $c=\cardinalSequence{\kappa}{d}$.
    Let $q=\cardinalSequenceEnvelope{\kappa}{d}$.
    Let $i\in\lambda$.
    Suppose $i$ is not a limit ordinal.
    Then $q(i)\in c(i)$.
\end{proposition}
\begin{proof}
    $q\in\funs{\lambda}{\kappa}$
        by \auto{cardinal_sequence_envelope_type}.
    Let $g=\restrl{c}{i}$.
    $g\in\funs{i}{\kappa}$
        by \auto{cardinal_sequence_restriction_type}.
    Let $h=\cardinalSequenceNext{q}{\lambda}{i}{g}$.
    $q(i)\in h$ by \auto{cardinal_sequence_next_properties}.
    $c(i)=h$ by \auto{cardinal_sequence_at_nonlimit}.
\end{proof}

\begin{proposition}\label{cardinal_sequence_is_cofinal}
    Let $\kappa$ be a singular cardinal.
    Let $\lambda=\cofinality{\kappa}$.
    Suppose $\cofinalMap{d}{\lambda}{\kappa}$.
    Let $c=\cardinalSequence{\kappa}{d}$.
    Then $\cofinalMap{c}{\lambda}{\kappa}$.
\end{proposition}
\begin{proof}
    $c\in\funs{\lambda}{\kappa}$ by \auto{cardinal_sequence_type}.
    Hence $c$ is a function and $\dom{c}=\lambda$
        by \obv{funs_is_function,funs}.
    $\ran{c}\subseteq\kappa$ by \obv{funs_ran}.
    Hence $\kappa$ is a limit ordinal
        by \obv{infinite_cardinal_is_limit_ordinal}.
    $\lambda$ is an infinite cardinal
        by \obv{cofinality_of_limit_is_infinite_cardinal}.
    Hence $\lambda$ is a limit ordinal
        by \obv{infinite_cardinal_is_limit_ordinal}.
    Show that for every $\alpha\in\kappa$ there exists $\gamma\in\ran{c}$
        such that $\alpha\subseteq\gamma$.
    \begin{subproof}
        Fix $\alpha\in\kappa$.
        $\cofinalIn{\ran{d}}{\kappa}$ by \obv{cofinal_map}.
        Take $j\in\ran{d}$ such that $\alpha\subseteq j$
            by \obv{cofinal_in}.
        $d\in\funs{\lambda}{\kappa}$ by \obv{cofinal_map}.
        Hence $d$ is a function and $\dom{d}=\lambda$
            by \obv{funs_is_function,funs}.
        Take $r\in\lambda$ such that $d(r)=j$
            by \auto{fun_ran_iff}.
        Take $i\in\lambda$ such that $i$ is not a limit ordinal and
            $r\in\suc{i}$ by \obv{nonlimit_index_above}.
        $d(r)\in(\cardinalSequenceEnvelope{\kappa}{d})(i)$
            by \obv{cardinal_sequence_envelope_dominates}.
        $(\cardinalSequenceEnvelope{\kappa}{d})(i)\in c(i)$
            by \obv{cardinal_sequence_dominates_envelope_at_nonlimit}.
        Hence $d(r)\in c(i)$
            by \auto{ordinal_transitivity,cardinal_sequence_stage_properties,cardinal_sequence_stage_property,cardinal,initial_ordinal}.
        Hence $j\in c(i)$.
        $c(i)$ is an ordinal
            by \auto{cardinal_sequence_stage_properties,cardinal_sequence_stage_property,cardinal,initial_ordinal}.
        Hence $j\subseteq c(i)$
            by \obv{elem_of_ordinal_is_ordinal,ordinal_elem_implies_subseteq}.
        Hence $\alpha\subseteq c(i)$ by \obv{subseteq_transitive}.
        $c(i)\in\ran{c}$ by \auto{fun_ran_iff}.
        Follows by \auto.
    \end{subproof}
    Hence $\cofinalIn{\ran{c}}{\kappa}$ by \obv{cofinal_in}.
    Follows by \obv{cofinal_map}.
\end{proof}

\begin{proposition}\label{cardinal_sequence_is_index_dominating}
    Let $\kappa$ be a singular cardinal.
    Let $\lambda=\cofinality{\kappa}$.
    Suppose $\cofinalMap{d}{\lambda}{\kappa}$.
    Let $c=\cardinalSequence{\kappa}{d}$.
    Then $c$ is index-dominating on $\lambda$.
\end{proposition}
\begin{proof}
    Hence $\kappa$ is a limit ordinal
        by \obv{infinite_cardinal_is_limit_ordinal}.
    Hence $\lambda$ is an infinite cardinal
        by \obv{cofinality_of_limit_is_infinite_cardinal}.
    Hence $\lambda$ is a limit ordinal
        by \obv{infinite_cardinal_is_limit_ordinal}.
    Hence $\lambda$ is an ordinal by \obv{limit_ordinal}.
    $\lambda\in c(\emptyset)$
        by \obv{cardinal_sequence_contains_cofinality_at_zero}.
    Show that for every $i\in\lambda$ we have $\card{i}\in c(i)$.
    \begin{subproof}
        Fix $i\in\lambda$.
        Hence $i$ is an ordinal by \obv{elem_of_ordinal_is_ordinal}.
        $\card{i}\in\lambda$ by \obv{card_of_ordinal_below}.
        Show that $\lambda\in c(i)$.
        \begin{subproof}
            \begin{byCase}
                \caseOf{$i=\emptyset$.}
                    Follows.
                \caseOf{$i\neq\emptyset$.}
                    Hence $\emptyset\in i$
                        by \obv{ordinal_empty_or_emptyset_elem}.
                    Hence $c(\emptyset)\in c(i)$
                        by \auto{cardinal_sequence_stage_properties,cardinal_sequence_stage_property}.
                    $c(i)$ is an infinite cardinal
                        by \auto{cardinal_sequence_stage_properties,cardinal_sequence_stage_property}.
                    Hence $c(i)$ is an ordinal
                        by \auto{cardinal,initial_ordinal}.
                    Follows by \auto{ordinal_transitivity}.
            \end{byCase}
        \end{subproof}
        $c(i)$ is an infinite cardinal
            by \auto{cardinal_sequence_stage_properties,cardinal_sequence_stage_property}.
        Hence $c(i)$ is an ordinal by \auto{cardinal,initial_ordinal}.
        Follows by \auto{ordinal_transitivity}.
    \end{subproof}
    Follows by \obv{index_dominating}.
\end{proof}

\begin{theorem}\label{cardinal_sequence_is_normal_cofinal}
    Let $\kappa$ be a singular cardinal.
    Let $\lambda=\cofinality{\kappa}$.
    Suppose $\cofinalMap{d}{\lambda}{\kappa}$.
    Let $c=\cardinalSequence{\kappa}{d}$.
    Then $\normalCofinalCardinalSequence{c}{\lambda}{\kappa}$ and
        $c$ is index-dominating on $\lambda$ and
        $\infiniteCardinalSequenceValues{c}{\lambda}$.
\end{theorem}
\begin{proof}
    $c\in\funs{\lambda}{\kappa}$ by \auto{cardinal_sequence_type}.
    $\lambda$ is a cardinal by \auto{cofinality_is_cardinal}.
    Hence $\lambda$ is an ordinal by \auto{cardinal,initial_ordinal}.
    Hence $\kappa$ is an ordinal by \auto{cardinal,initial_ordinal}.
    Hence $\ordinalMap{c}{\lambda}{\kappa}$ by \obv{ordinal_map}.
    $\cofinalMap{c}{\lambda}{\kappa}$
        by \obv{cardinal_sequence_is_cofinal}.
    For every $i\in\lambda$ we have $c(i)$ is an infinite cardinal
        by \auto{cardinal_sequence_stage_properties,cardinal_sequence_stage_property}.
    Hence $\infiniteCardinalSequenceValues{c}{\lambda}$
        by \obv{infinite_cardinal_sequence_values}.
    Hence $\cardinalSequenceValues{c}{\lambda}$
        by \auto{cardinal_sequence_values}.
    For all $j,k$ such that $j\in k\in\lambda$ we have $c(j)\in c(k)$
        by \auto{cardinal_sequence_stage_properties,cardinal_sequence_stage_property}.
    Hence $\normalMapIncreasing{c}{\lambda}$
        by \obv{normal_map_increasing}.
    For every limit ordinal $l\in\lambda$ we have
        $c(l)=\unions{\ran{\restrl{c}{l}}}$
        by \obv{cardinal_sequence_at_limit}.
    Hence $\normalMapContinuous{c}{\lambda}$
        by \obv{normal_map_continuous}.
    Hence $\normalMap{c}{\lambda}{\kappa}$ by \obv{normal_map}.
    Hence $\normalCofinalMap{c}{\lambda}{\kappa}$
        by \obv{normal_cofinal_map}.
    Hence $\normalCofinalCardinalSequence{c}{\lambda}{\kappa}$
        by \obv{normal_cofinal_cardinal_sequence}.
    $c$ is index-dominating on $\lambda$
        by \obv{cardinal_sequence_is_index_dominating}.
\end{proof}

\begin{theorem}\label{normal_cofinal_cardinal_sequence_exists}
    Let $\kappa$ be a singular cardinal.
    Let $\lambda=\cofinality{\kappa}$.
    Then there exists $c$ such that
        $\normalCofinalCardinalSequence{c}{\lambda}{\kappa}$.
\end{theorem}
\begin{proof}
    Take $d$ such that $\cofinalMap{d}{\lambda}{\kappa}$
        by \auto{cofinality_map_exists}.
    Let $c=\cardinalSequence{\kappa}{d}$.
    $\normalCofinalCardinalSequence{c}{\lambda}{\kappa}$
        by \auto{cardinal_sequence_is_normal_cofinal}.
    Follows by \auto.
\end{proof}

\begin{theorem}\label{index_dominating_normal_cofinal_cardinal_sequence_exists}
    Let $\kappa$ be a singular cardinal.
    Let $\lambda=\cofinality{\kappa}$.
    Then there exists $c$ such that
        $\normalCofinalCardinalSequence{c}{\lambda}{\kappa}$ and
        $c$ is index-dominating on $\lambda$ and
        $\infiniteCardinalSequenceValues{c}{\lambda}$.
\end{theorem}
\begin{proof}
    Take $d$ such that $\cofinalMap{d}{\lambda}{\kappa}$
        by \auto{cofinality_map_exists}.
    Let $c=\cardinalSequence{\kappa}{d}$.
    $\normalCofinalCardinalSequence{c}{\lambda}{\kappa}$ and
        $c$ is index-dominating on $\lambda$ and
        $\infiniteCardinalSequenceValues{c}{\lambda}$
        by \obv{cardinal_sequence_is_normal_cofinal}.
    Follows by \auto.
\end{proof}
